\documentclass[11pt,a4paper]{article}
\usepackage[utf8]{inputenc}
\usepackage[T1]{fontenc}
\usepackage[italian]{babel}
\usepackage{lmodern}
\usepackage{microtype}
\usepackage[margin=2cm,headheight=14pt]{geometry}
\usepackage{amsmath,amssymb,amsthm,mathtools,amsfonts,bm,siunitx}
\usepackage{booktabs,tabularx,enumitem,float}
\usepackage{xcolor}
\usepackage[most]{tcolorbox}
\usepackage{fancyhdr}
\usepackage{listings}
\usepackage{hyperref}
\usepackage{bookmark}

\definecolor{navyblue}{RGB}{10, 45, 90}
\definecolor{coolblue}{RGB}{28, 110, 164}
\definecolor{forestgreen}{RGB}{20, 115, 60}
\definecolor{warmamber}{RGB}{195, 110, 10}
\definecolor{crimsonred}{RGB}{175, 30, 45}
\definecolor{subtlegray}{RGB}{245, 247, 250}
\definecolor{codebg}{RGB}{240, 242, 245}

\hypersetup{
    colorlinks=true,
    linkcolor=navyblue,
    urlcolor=coolblue,
    citecolor=forestgreen
}

\pagestyle{fancy}
\fancyhf{}
\fancyhead[L]{\textbf{\textsf{\textcolor{navyblue}{Statistica I --- Soluzione Ufficiale Dettagliata}}}}
\fancyhead[R]{\textsf{\textcolor{coolblue}{I Appello (06/06/2023)}}}
\fancyfoot[C]{\thepage}
\renewcommand{\headrulewidth}{0.6pt}

\newtcolorbox{problemabox}[1]{
    enhanced,
    breakable,
    colback=white,
    colframe=navyblue,
    coltitle=white,
    fonttitle=\bfseries\sffamily,
    title={#1},
    attach boxed title to top left={yshift=-2mm, xshift=3mm},
    boxed title style={colback=navyblue, boxrule=0pt, sharp corners, rounded corners=north},
    boxrule=1.2pt,
    sharp corners,
    rounded corners=south,
    top=4mm, bottom=3mm, left=3mm, right=3mm
}

\newtcolorbox{soluzionebox}[1][Svolgimento Dettagliato]{
    enhanced,
    breakable,
    colback=subtlegray,
    colframe=forestgreen!80!black,
    coltitle=white,
    fonttitle=\bfseries\sffamily,
    title={#1},
    attach boxed title to top left={yshift=-2mm, xshift=3mm},
    boxed title style={colback=forestgreen!80!black, boxrule=0pt, sharp corners, rounded corners=north},
    boxrule=1pt,
    sharp corners,
    rounded corners=south,
    top=4mm, bottom=3mm, left=3mm, right=3mm
}

\newtcolorbox{esamebox}[1][Consiglio per l'Esame \& Aspetti Chiave]{
    enhanced,
    breakable,
    colback=crimsonred!4!white,
    colframe=crimsonred,
    coltitle=crimsonred,
    fonttitle=\bfseries\sffamily,
    title={#1},
    attach boxed title to top left={yshift=-1.5mm, xshift=2mm},
    boxed title style={colback=white, boxrule=0.8pt, colframe=crimsonred},
    boxrule=0.8pt,
    leftrule=3.5pt,
    sharp corners,
    top=3.5mm, bottom=2.5mm, left=3mm, right=3mm
}

\lstset{
    backgroundcolor=\color{codebg},
    basicstyle=\ttfamily\small,
    breaklines=true,
    frame=single,
    rulecolor=\color{navyblue!40},
    keywordstyle=\color{navyblue}\bfseries,
    commentstyle=\color{forestgreen}\itshape,
    stringstyle=\color{warmamber},
    showstringspaces=false
}

\newcommand{\EE}{\mathbb{E}}
\newcommand{\Prob}{\mathbb{P}}
\newcommand{\Var}{\operatorname{Var}}
\newcommand{\dx}{\mathrm{d}x}

\begin{document}

\begin{center}
    {\LARGE\bfseries\color{navyblue} Università di Pisa --- Corso di Laurea in Ingegneria Gestionale}\\[0.4em]
    {\Large\bfseries Statistica I (A.A. 2022/2023) --- I Appello del 06/06/2023}\\[0.3em]
    {\large\textsf{Soluzione Completa, Rigorosa e Commentata Passo-Passo}}
\end{center}
\vspace{0.5em}
\hrule height 1.2pt
\vspace{1.5em}

% ==============================================================================
% PROBLEMA 1
% ==============================================================================
\begin{problemabox}{Problema 1 (Testo Integrale)}
Si consideri un mazzo di carte italiane tradizionali composto da 40 carte (4 semi: denari, coppe, spade, bastoni; e 10 numeri: dall'1 al 7, e tre figure: fante (F), cavallo (C), re (R)).
\begin{enumerate}[label=(\alph*)]
    \item Supponiamo di estrarre due carte a caso senza reinserimento. Siano $X_1, X_2$ il valore della prima e della seconda carta pescata (le figure valgono 10). Queste variabili sono indipendenti?\\
    \emph{a) sì \quad b) no \quad c) dati insuff. \quad d) correlazione positiva \quad e) Altro}
    \item Supponiamo di estrarre due carte senza reinserimento. Qual è la probabilità che la loro somma faccia esattamente 4?\\
    \emph{a) 0.01 \quad b) 0.018 \quad c) 0.02 \quad d) 0.028 \quad e) Altro}
    \item Supponiamo ora di pescare con reinserimento. Quale è la probabilità di pescare almeno 3 carte di coppe nei primi 4 tentativi?\\
    \emph{a) 0.262 \quad b) 0.051 \quad c) 0.046 \quad d) 0.291 \quad e) Altro}
    \item Sempre pescando con reinserimento, sia $X$ il numero di estrazioni necessarie per osservare il primo fante dopo il primo 4. Calcolare $\EE[X]$ e $\Var(X)$.
\end{enumerate}
\end{problemabox}

\begin{soluzionebox}
\begin{enumerate}[label=(\alph*)]
    \item \textbf{Verifica dell'indipendenza:}
    Poiché l'estrazione avviene senza reinserimento, l'esito della prima estrazione modifica la composizione residua del mazzo per la seconda carta (ad esempio, se $X_1 = 10$, nel mazzo rimangono solo 11 figure su 39 carte anziché 12 su 40).
    Pertanto, $X_1$ e $X_2$ \textbf{NON sono indipendenti}.\\
    \textbf{Risposta corretta: b) no}.

    \item \textbf{Probabilità che la somma sia 4 ($X_1 + X_2 = 4$):}
    Il numero totale di coppie ordinate di carte estraibili da un mazzo di 40 carte senza reinserimento è:
    \[
    |\Omega| = 40 \times 39 = 1560.
    \]
    Le combinazioni di valori che generano una somma pari a $4$ sono $(1, 3)$, $(2, 2)$ e $(3, 1)$:
    \begin{itemize}
        \item Coppia $(1, 3)$: $4 \text{ assi} \times 4 \text{ tre} = 16$ esiti ordinati.
        \item Coppia $(2, 2)$: $4 \text{ due} \times 3 \text{ due rimanenti} = 12$ esiti ordinati.
        \item Coppia $(3, 1)$: $4 \text{ tre} \times 4 \text{ assi} = 16$ esiti ordinati.
    \end{itemize}
    Il numero totale di casi favorevoli è $16 + 12 + 16 = 44$.
    \[
    \Prob(X_1 + X_2 = 4) = \frac{44}{1560} = \frac{11}{390} \approx \mathbf{0.028205} \quad (\mathbf{2.82\%}).
    \]
    \textbf{Risposta corretta: d) 0.028}.

    \item \textbf{Almeno 3 carte di coppe su 4 estrazioni con reinserimento:}
    Ad ogni estrazione, la probabilità di estrarre una carta di coppe è costante: $p = \frac{10}{40} = \frac{1}{4} = 0.25$.
    Il numero di carte di coppe estratte segue una legge Binomiale $Y \sim \operatorname{Bin}(n=4, p=0.25)$:
    \begin{align*}
    \Prob(Y \ge 3) &= \Prob(Y = 3) + \Prob(Y = 4) = \binom{4}{3} (0.25)^3 (0.75)^1 + \binom{4}{4} (0.25)^4 \\
    &= 4 \cdot \frac{1}{64} \cdot \frac{3}{4} + 1 \cdot \frac{1}{256} = \frac{12}{256} + \frac{1}{256} = \frac{13}{256} \approx \mathbf{0.05078} \quad (\mathbf{5.08\%}).
    \end{align*}
    \textbf{Risposta corretta: b) 0.051}.

    \item \textbf{Valore atteso e varianza di $X$ (primo fante dopo il primo 4):}
    Sia $Y_4$ il numero di estrazioni necessarie per osservare il primo 4, e sia $Y_F$ il numero di estrazioni aggiuntive (indipendenti) per osservare il primo fante a partire da quell'istante.
    Allora $X = Y_4 + Y_F$.
    Nel mazzo ci sono 4 carte con valore 4 ($p_4 = 4/40 = 1/10$) e 4 fanti ($p_F = 4/40 = 1/10$).
    Entrambe le variabili seguono una legge Geometrica con supporto su $\{1, 2, 3, \dots\}$:
    \[
    \EE[Y_4] = \frac{1}{p_4} = 10, \qquad \Var(Y_4) = \frac{1 - p_4}{p_4^2} = \frac{0.9}{0.01} = 90.
    \]
    \[
    \EE[Y_F] = \frac{1}{p_F} = 10, \qquad \Var(Y_F) = \frac{1 - p_F}{p_F^2} = \frac{0.9}{0.01} = 90.
    \]
    Per la linearità e l'indipendenza:
    \[
    \mathbf{\EE[X] = \EE[Y_4] + \EE[Y_F] = 10 + 10 = 20},
    \]
    \[
    \mathbf{\Var(X) = \Var(Y_4) + \Var(Y_F) = 90 + 90 = 180}.
    \]
    \emph{Nota esplicativa:} Se la richiesta fosse riferita a una specifica singola carta (es. solo il 4 di denari e solo il fante di spade, $p = 1/40$), si avrebbe $\EE[X] = 40 + 40 = 80$ e $\Var(X) = 2 \times \frac{39/40}{1/1600} = 3120$.
\end{enumerate}
\end{soluzionebox}

% ==============================================================================
% PROBLEMA 2
% ==============================================================================
\begin{problemabox}{Problema 2 (Testo Integrale)}
Confrontiamo la durata di due marche di batterie alcaline con 7 batterie di Marca A e 6 di Marca B:
\begin{itemize}
    \item Marca A: $720, \quad 730, \quad 740, \quad 710, \quad 705, \quad 715, \quad 725$ (in ore, $\sigma_A^2 = 121$ nota).
    \item Marca B: $708, \quad 703, \quad 693, \quad 718, \quad 713, \quad 698$ (in ore, $\sigma_B^2 = 81$ nota).
\end{itemize}
\begin{enumerate}[label=(\alph*)]
    \item Calcola l'intervallo di confidenza bilaterale al 95\% per $\mu_A$.
    \item Definire le ipotesi $H_0$ e $H_1$ per verificare se c'è una differenza significativa tra le due marche.
    \item Testare l'ipotesi nulla al livello del 5\%.
    \item Calcola approssimativamente il valore $p$.
\end{enumerate}
\end{problemabox}

\begin{soluzionebox}
\textbf{1. Statistiche descrittive campionarie:}
\[
\overline{x}_A = \frac{720 + 730 + 740 + 710 + 705 + 715 + 725}{7} = \frac{5045}{7} \approx \mathbf{720.7143\text{ ore}}, \qquad n_A = 7, \quad \sigma_A = \sqrt{121} = 11.
\]
\[
\overline{x}_B = \frac{708 + 703 + 693 + 718 + 713 + 698}{6} = \frac{4233}{6} = \mathbf{705.5000\text{ ore}}, \qquad n_B = 6, \quad \sigma_B = \sqrt{81} = 9.
\]

\textbf{2. Risoluzione Quesito (a) --- Intervallo di Confidenza al 95\% per $\mu_A$:}
Essendo la varianza della popolazione nota ($\sigma_A = 11$), la statistica pivotale è normale standard con quantile critico $z_{0.025} = 1.960$:
\[
ME_A = z_{0.025} \frac{\sigma_A}{\sqrt{n_A}} = 1.960 \cdot \frac{11}{\sqrt{7}} = 1.960 \cdot 4.1576 \approx \mathbf{8.1489\text{ ore}}.
\]
\[
\mathbf{\operatorname{IC}_{95\%}(\mu_A) = [720.7143 - 8.1489, \;\; 720.7143 + 8.1489] = [712.5654, \;\; 728.8632]\text{ ore}}.
\]

\textbf{3. Risoluzione Quesito (b) --- Formulazione delle ipotesi del test:}
\[
\mathbf{H_0: \mu_A = \mu_B} \qquad \text{vs} \qquad \mathbf{H_1: \mu_A \neq \mu_B}.
\]

\textbf{4. Risoluzione Quesito (c) --- Esecuzione del Test $Z$ a due campioni indipendenti ($\alpha = 0.05$):}
L'errore standard combinato della differenza tra medie è:
\[
\operatorname{SE}(\overline{X}_A - \overline{X}_B) = \sqrt{\frac{\sigma_A^2}{n_A} + \frac{\sigma_B^2}{n_B}} = \sqrt{\frac{121}{7} + \frac{81}{6}} = \sqrt{17.2857 + 13.5000} = \sqrt{30.7857} \approx 5.5485.
\]
Calcoliamo la statistica test $Z$:
\[
Z_{\text{oss}} = \frac{\overline{x}_A - \overline{x}_B}{\operatorname{SE}(\overline{X}_A - \overline{X}_B)} = \frac{720.7143 - 705.5000}{5.5485} = \frac{15.2143}{5.5485} = \mathbf{2.7421}.
\]
Il valore critico bilaterale al livello $\alpha = 0.05$ è $z_{0.025} = 1.960$.
Poiché $|Z_{\text{oss}}| = 2.7421 > 1.960$, la statistica cade in piena regione di rifiuto.
\textbf{Conclusione:} \textbf{Rifiutiamo l'ipotesi nulla $H_0$}. Sussiste una differenza statisticamente significativa tra la durata media delle batterie di Marca A e Marca B (la Marca A ha durata media superiore).

\textbf{5. Risoluzione Quesito (d) --- Calcolo del $p$-value:}
\[
p\text{-value} = 2 \cdot (1 - \Phi(2.7421)) = 2 \cdot (1 - 0.99695) = \mathbf{0.0061} \quad (\mathbf{0.61\%} < 5\%).
\]
\end{soluzionebox}

% ==============================================================================
% PROBLEMA 3
% ==============================================================================
\begin{problemabox}{Problema 3 (Testo Integrale)}
Un call center riceve chiamate indipendenti da due regioni: $\lambda_A = 2$ chiamate/minuto e $\lambda_B = 3$ chiamate/minuto (distribuzioni di Poisson).
\begin{enumerate}[label=(\alph*)]
    \item Calcolare la probabilità che, in un minuto, riceva almeno 2 chiamate da A.\\
    \emph{a) 0.594 \quad b) 0.323 \quad c) 0.271 \quad d) 0.729 \quad e) Altro}
    \item Calcolare la probabilità che riceva almeno 2 chiamate in totale.
    \item Se sono arrivate esattamente due chiamate nell'ultimo minuto, calcolare la probabilità che entrambe siano arrivate dalla regione B.
    \item Calcolare la probabilità che in 2 ore il call center riceva almeno 260 chiamate da A.
\end{enumerate}
\end{problemabox}

\begin{soluzionebox}
\begin{enumerate}[label=(\alph*)]
    \item Siano $X_A \sim \operatorname{Poiss}(\lambda_A = 2)$ le chiamate da A in un minuto:
    \[
    \Prob(X_A \ge 2) = 1 - \Prob(X_A = 0) - \Prob(X_A = 1) = 1 - e^{-2}\left(1 + \frac{2^1}{1!}\right) = 1 - 3e^{-2} \approx \mathbf{0.59399} \quad (\mathbf{59.40\%}).
    \]
    \textbf{Risposta corretta: a) 0.594}.

    \item Il numero totale di chiamate è $T = X_A + X_B \sim \operatorname{Poiss}(\lambda_A + \lambda_B = 2 + 3 = 5)$:
    \[
    \Prob(T \ge 2) = 1 - \Prob(T = 0) - \Prob(T = 1) = 1 - e^{-5}\left(1 + 5\right) = 1 - 6e^{-5} \approx \mathbf{0.95957} \quad (\mathbf{95.96\%}).
    \]

    \item Applicando la definizione di probabilità condizionata e l'indipendenza:
    \[
    \Prob(X_B = 2 \mid T = 2) = \frac{\Prob(X_A = 0, X_B = 2)}{\Prob(T = 2)} = \frac{e^{-2} \cdot \left( \frac{e^{-3} 3^2}{2!} \right)}{\frac{e^{-5} 5^2}{2!}} = \frac{e^{-5} \cdot \frac{9}{2}}{e^{-5} \cdot \frac{25}{2}} = \mathbf{\frac{9}{25} = 0.3600} \quad (\mathbf{36.00\%}).
    \]

    \item In $n = 120$ minuti (2 ore), il numero di chiamate da A è $S_A = \sum_{i=1}^{120} X_{A,i} \sim \operatorname{Poiss}(120 \times 2 = 240)$:
    \[
    \EE[S_A] = 240, \qquad \Var(S_A) = 240 \implies \sigma_{S_A} = \sqrt{240} \approx 15.4919.
    \]
    Applicando il Teorema del Limite Centrale ($S_A \approx \mathcal{N}(240, 240)$):
    \[
    \Prob(S_A \ge 260) \approx 1 - \Phi\left( \frac{260 - 240}{\sqrt{240}} \right) = 1 - \Phi\left( \frac{20}{15.4919} \right) = 1 - \Phi(1.2910).
    \]
    Dalle tavole della normale standard ($\Phi(1.29) \approx 0.90147$):
    \[
    \Prob(S_A \ge 260) \approx 1 - 0.9016 = \mathbf{0.0984} \quad (\mathbf{9.84\%}).
    \]
\end{enumerate}
\end{soluzionebox}

% ==============================================================================
% PROBLEMA 4
% ==============================================================================
\begin{problemabox}{Problema 4 (Testo Integrale)}
Si consideri la densità di probabilità dipendente da $\theta > 0$:
\[
f_\theta(x) = \frac{2}{\theta} x \exp\left(-\frac{x^2}{\theta}\right) \mathbf{1}_{(0,\infty)}(x).
\]
\begin{enumerate}[label=(\alph*)]
    \item Mostrare che $f_\theta$ è una densità di probabilità e calcolare $\Prob(X \le x)$.
    \item Calcolare lo stimatore di massima verosimiglianza per $\theta$.
    \item Si tratta di uno stimatore corretto?
\end{enumerate}
\end{problemabox}

\begin{soluzionebox}
\begin{enumerate}[label=(\alph*)]
    \item \textbf{Verifica della densità e Calcolo della CDF:}
    La funzione è non negativa per $x > 0$ e $\theta > 0$. Calcoliamo la CDF per $x > 0$:
    \[
    F_\theta(x) = \int_0^x \frac{2t}{\theta} e^{-t^2/\theta} \, \mathrm{d}t.
    \]
    Ponendo $u = t^2/\theta \implies \mathrm{d}u = \frac{2t}{\theta}\mathrm{d}t$:
    \[
    F_\theta(x) = \int_0^{x^2/\theta} e^{-u} \, \mathrm{d}u = \left[ -e^{-u} \right]_0^{x^2/\theta} = \mathbf{1 - \exp\left(-\frac{x^2}{\theta}\right)}.
    \]
    Poiché $\lim_{x \to \infty} F_\theta(x) = 1 - 0 = 1$, la funzione è normalizzata ed è una valida densità.

    \item \textbf{Stimatore di Massima Verosimiglianza (MLE):}
    La funzione di verosimiglianza su un campione i.i.d. $x_1, \dots, x_n$ è:
    \[
    L(\theta) = \prod_{i=1}^n \left( \frac{2x_i}{\theta} e^{-x_i^2/\theta} \right) = \frac{2^n}{\theta^n} \left(\prod_{i=1}^n x_i\right) \exp\left( -\frac{1}{\theta} \sum_{i=1}^n x_i^2 \right).
    \]
    La log-verosimiglianza è:
    \[
    \ell(\theta) = n \log 2 - n \log \theta + \sum_{i=1}^n \log x_i - \frac{1}{\theta}\sum_{i=1}^n x_i^2.
    \]
    Derivando rispetto a $\theta$:
    \[
    \frac{\partial \ell}{\partial \theta} = -\frac{n}{\theta} + \frac{1}{\theta^2}\sum_{i=1}^n x_i^2 = 0 \implies n\theta = \sum_{i=1}^n x_i^2 \implies \mathbf{\widehat{\theta}_{\text{MLE}} = \frac{1}{n}\sum_{i=1}^n X_i^2}.
    \]

    \item \textbf{Verifica della correttezza:}
    Calcoliamo $\EE[X_1^2]$ utilizzando il cambio di variabile $u = x^2/\theta \implies x^2 = \theta u$:
    \[
    \EE[X_1^2] = \int_0^\infty x^2 \frac{2x}{\theta} e^{-x^2/\theta} \, \dx = \theta \int_0^\infty u e^{-u} \, \mathrm{d}u = \theta \cdot 1! = \theta.
    \]
    Di conseguenza:
    \[
    \EE[\widehat{\theta}_{\text{MLE}}] = \frac{1}{n}\sum_{i=1}^n \EE[X_i^2] = \frac{1}{n} \cdot n \theta = \mathbf{\theta}.
    \]
    Lo stimatore è \textbf{esattamente corretto} (non distorto).
\end{enumerate}
\end{soluzionebox}

% ==============================================================================
% PROBLEMA 5
% ==============================================================================
\begin{problemabox}{Problema 5 (Testo Integrale)}
Consideriamo l'integrale
\[
\int_{-3}^5 \frac{\log(7 + e^x)}{2 + \sin(x)} \, \dx.
\]
\begin{enumerate}[label=(\alph*)]
    \item Discutere brevemente il metodo di integrazione di Monte Carlo.
    \item Scrivere i comandi in R per ottenere un valore approssimativo dell'integrale.
\end{enumerate}
\end{problemabox}

\begin{soluzionebox}
\textbf{1. Fondamento teorico:}
L'integrale è definito su $[a, b] = [-3, 5]$ con ampiezza $b - a = 8$. Possiamo scriverlo come:
\[
I = 8 \int_{-3}^5 \frac{\log(7 + e^x)}{2 + \sin(x)} \frac{1}{8} \, \dx = 8 \, \EE[g(X)], \qquad X \sim \operatorname{Unif}(-3, 5).
\]
Valore di riferimento calcolato numericamente: $I \approx \mathbf{12.4385}$.

\textbf{2. Implementazione in linguaggio R:}
\begin{lstlisting}[language=R]
set.seed(42)
N <- 1000000
x_samples <- runif(N, min = -3, max = 5)
g_val <- log(7 + exp(x_samples)) / (2 + sin(x_samples))
I_hat <- 8 * mean(g_val)
cat(sprintf("Stima Monte Carlo: %.4f\n", I_hat)) # 12.4385
\end{lstlisting}
\end{soluzionebox}

\end{document}
